\documentclass[12pt]{article}
\usepackage{float}
\usepackage{graphicx}
\usepackage{fontspec}
\usepackage{url}
\usepackage{apacite}
\setmainfont{Times New Roman}
\usepackage[letterpaper,top=2cm,bottom=2cm,left=2cm,right=2cm,marginparwidth=1.5cm]{geometry}

\title{Weekly assignment 9: Ethics and Surveillance}
\author{Truong Khang Tobias Nguyen}
\date{}

\begin{document}

\maketitle

\section{Acceptable Surveillance}
It is difficult to say when surveillance is acceptable or not from a holistic view. On one hand, surveillance is a common part of our lives as governments and companies use surveillance technologies to monitor activities, collect information, and identify potential threats. In that sense, CCTVs, location trackers, and network monitoring systems have become common applications and I do not believe that we could live normally without these.
\newline
On the other hand, even though these kinds of surveillance technologies are employed with security and prevention in mind, they can also impede privacy and personal freedom and therefore, the question is not simply whether when surveillance is acceptable, but under which circumstances it is necessary.
\newline
I believe that surveillance is ethical when it serves a legitimate purpose and is done in a transparent, proportionate, and responsible manner. For example, CCTVs in public places are typically employed to help prevent crime or identify offenders which I consider a legitimate usage of surveillance. Another example, organizations that monitor their networks to identify attacks or suspicious activities is also a legitimate reason. Examples like these demonstrate that surveillance can be used to protect individuals or society as long as it is conducted with their best interest in mind.
\newline
However, like I have said in earlier assignments, this kind of power also comes with a lot of responsibility. Even though surveillance can be used for protection, it can also be used to monitor people and ultimately restrict their freedom. That being said, I believe that this question ultimately boils down to proportionality - the amount of surveillance should correspond to the actual risk or threat. For example, network monitoring to identify suspicious activity may be justified, while reading every person's private messages without reason is excessive and borders to abuse of power in my opinion. That being said, I believe that surveillance is acceptable as long as it is justified in it's reasoning and is limited to information that is actually relevant.
\newline
But I think it is also important that people should be aware if they are being monitored, what information or activities in particular are monitored and why they are being monitored in the first place. On that note, I consider surveillance unacceptable/unethical when it used as a tool of manipulation, intimidation, or excessive control. This may create a feeling of constant monitoring, which may force people to adapt their behaviour out of fear of repercussions, ultimately affecting their freedom.

\section{Surveillance in Cybersecurity}
To continue on the previous section, I see surveillance as a necessity in cybersecurity. Organizations not only have a responsibility to protect their own assets but these assets are often bound to other people's information that not everyone should access. To protect assets and information, organizations have to employ measures such as network monitoring to identify suspicious behaviour, or enable logging to trace incidents. 
\newline
But the same thing I stated in the previous section also applies in the context of cybersecurity. Surveillance in the name of cybersecurity is acceptable/ethical when it serves a legitimate purpose. As organizations have a high interest in protecting their assets, they are justified in employing surveillance technologies to ensure this protection as any compromise in them will affect individual rights.
\newline
Again, it is \textbf{essential} that any kind of surveillance technology is clearly communicated to the people that are "targeted" by it and that they are aware what actions are monitored, what information is necessary and why it matters. Additionally, it is also important that those in power of the surveillance are aware of their duties and fulfill them in a responsible manner, that they only use such technologies for their intended purposes and that they are aware that there are consequences with they abuse this power.
\newline
The moment surveillance is used for other purposes that may impede an individual's freedom such as tracking productivity, monitoring interactions \& conversations, or bulk data collection, it becomes an abuse of power and thus, ends the acceptance of surveillance.

\section{Contextual Influences}
Our environment itself influences the acceptance of surveillance as different contexts have different purposes, proportionalities, and requirements in regards to surveillance.
\newline
In a workplace environment, the organization should use surveillance technologies in order to protect assets. It is normal for organizations to employ measures such as network monitoring, log collection or threat intelligence. In this context, surveillance becomes a "necessary evil" as it is needed to protect these assets and the consequences of a compromise may be greater than using the surveillance itself.
\newline
In a home environment, the requirements shift from protecting assets to protecting the inhabitants of a given home and the technology is typically found in the form of home cameras or sensors around the perimeters. In this context, the inhabitants don't use surveillance on themselves but on anyone who may threaten this home. Even though it threatens an adversaires individual rights, I argue that the moment one is ready to harm another person's rights, they frankly do not deserve protection themself.
\newline
In a societal context, surveillance becomes a trade-off between collective security vs. individual rights. While measures such as CCTVs or social media monitoring are used to identify prevent crimes or identify offenders, they may also be used for discrimination or disabling critical opinions which may cause innocent people to shift their own behaviour out of fear of judgement and thus, impedes an individual's right of expression. The trade-off becomes potentially sacrifing individual rights for "the greater good".
\newline
That being said, the acceptance of surveillance also depends on the context it is used. Even though moust scenarios differ in their purpose, proportionality, and requirements, the things I have stated in the first section still apply - surveillance is acceptable as long as it is justified in it's purpose and the people surveilled are aware of it.

\section{Regulatory and Organizational Influences}
The acceptability of surveillance is also influenced by regulatory and organizational factors. Lawmakers can decide what surveillance is allowed, who can conduct it, what information can be collected, and under what circumstances it becomes intrusive/abusive. We can already see how laws influence surveillance in many forms like privacy \& data protection laws such as the GDPR, workplace monitoring laws, or criminal laws - these laws constitute in which context surveillance is allowed to what extent and therefore, influence it.
\newline
Organizations, on the other hand, can influence surveillance as long as it remains lawful. They translate regulatory requirements into practical rules about surveillance and can adapt them to their specific context. They develop policies that determine what can be surveilled, who can conduct it, why they have to conduct, and how they inform individuals.
\newline
In that sense, laws and organizational policies shape what constitutes surveillance, it's boundaries, and purposes. While laws create a baseline, organizational policies determine how they are used in practice and adapt them to their own context.

\end{document} 